% Generated draft: requires mathematical and source review.
\documentclass[11pt]{article}
\usepackage[margin=25mm]{geometry}
\usepackage{fontspec}
\setmainfont{Latin Modern Roman}
\usepackage{amsmath,amssymb,mathtools,unicode-math}
\setmathfont{Latin Modern Math}
\usepackage{enumitem,needspace,xurl,hyperref}
\setlength{\parindent}{0pt}
\setlength{\parskip}{5pt}
\setlength{\emergencystretch}{3em}
\newcommand{\N}{\mathbb{N}}
\newcommand{\Z}{\mathbb{Z}}
\newcommand{\Q}{\mathbb{Q}}
\newcommand{\R}{\mathbb{R}}
\newcommand{\C}{\mathbb{C}}
\newcommand{\F}{\mathbb{F}}
\newcommand{\A}{\mathbb{A}}
\newcommand{\PP}{\mathbb{P}}
\newcommand{\E}{\mathbb{E}}
\newcommand{\eps}{\varepsilon}
\newcommand{\dd}{\,\mathrm{d}}
\begin{document}
\section*{UnsolvedMath ID: 11000306}
\textbf{Source problem code:} catalogue entry
\subsection{Definitions and mathematical statement}
Retain \(\mathfrak u_g\), described in the source as the graded Lie algebra of the prounipotent radical of the relative Malcev completion of \(\mathcal I_g\) defined by Hain [29]. Let \(\mathfrak h_g^\Q\) be the positively graded Lie algebra of symplectic derivations of the Malcev Lie algebra of \(\pi_1(\Sigma_g)\), with the source's natural homomorphism
\[
 \psi_g:\mathfrak u_g\longrightarrow\mathfrak h_g^\Q.
\]
\emph{Problem 3.3.} Determine whether \(\psi_g\) is injective, that is, whether \(\ker\psi_g=0\), in the intended genus range. The completion's reference representation, grading, and precise symplectic-derivation convention are not supplied. In particular, if \(\mathcal I_g\) means the Torelli group, its usual homology representation is trivial; relative completion of that group must not be confused with relative completion of the whole mapping class group. The printed group name must be checked rather than silently corrected.

The extraction then gives separate secondary-characteristic-class context. Write \(M_g\) for the mapping class group, \(\mathcal I_g\) for Torelli, \(U\) for the unspecified Johnson target, and \(\mathcal K_g\) for its kernel. The source describes a homomorphism \(M_g\to U\rtimes\operatorname{Sp}(2g,\Z)\) [77] extending the first Johnson map \(\mathcal I_g\to U\), and invariant cocycles
\[
 z_i\in Z^{2i}(M_g;\Q),\qquad [z_i]=e_i,
\]
constructed in [51,52] for the Mumford--Morita--Miller classes, said unique up to coboundaries. For odd-index classes it chooses invariant cocycles
\[
 z'_{2i-1}\in Z^{4i-2}(M_g;\Q)
\]
pulled back from \(\operatorname{Sp}(2g,\Z)\), representing the same classes as \(z_{2i-1}\). Thus it chooses
\[
 y_i\in C^{4i-3}(M_g;\Q),\qquad
 \delta y_i=z_{2i-1}-z'_{2i-1}.
\]
Its use of Madsen--Weiss [64] is expressly in a suitable stable range, where \(H^{4i-3}(M_g;\Q)=0\), so the stated ambiguity in \(y_i\) is a coboundary. No explicit stable genus bound is supplied.

It records the extension
\[
 1\longrightarrow\mathcal K_g\longrightarrow\mathcal I_g
 \longrightarrow U\longrightarrow1,\tag{8}
\]
and credits Johnson [43] with generation of \(\mathcal K_g\) by separating twists. With its chosen \(z_i\) vanishing on \(\mathcal K_g\) and \(z'_{2i-1}\) trivial on Torelli, restriction gives
\[
 d_i=[y_i|_{\mathcal K_g}]\in H^{4i-3}(\mathcal K_g;\Q).
\]
The source says \(d_i\) is \(M_g\)-invariant for the outer-conjugation action: \(\delta(\varphi^*y_i-y_i)=0\), and the same stable vanishing makes that difference a coboundary. It explains this as removing the ambiguity in the earlier secondary classes [80]. These are supplied contextual constructions with their cocycle choices and stable-range limitations; they do not answer the Lie-algebra injectivity problem.
\subsection{Short English statement}
Determine injectivity of the source's natural graded Lie-algebra map to positive-degree symplectic derivations. The completion setup needs checking. The adjacent secondary-class construction is retained with all degrees and its essential stable-range cohomology vanishing.
\subsection{Sources}
\begin{itemize}
\item Local UnsolvedMath registry, permanent ID 11000306, source code catalogue entry, Problem 3.3, PDF page 362. No individual proposer is supplied.
\item Farb, \emph{Problems on Mapping Class Groups and Related Topics}. Supplied URL: \url{https://www.math.uchicago.edu/~farb/papers/mcgbook.pdf}; recorded access 2026-07-29; extraction pdf-font-structure. The volume contains 329 explicitly typeset headings across chapters, with repeated labels disambiguated by page.
\item Hain [29] for the completion; [80] for earlier secondary classes; [51,52] for cocycles; [77] for the Johnson extension; Madsen--Weiss [64] for stable vanishing; Johnson [43] for separating-twist generation. Full bibliographic entries are absent.
\item Supplied status open and OPEN-TRIAGE. The summary describes a truncated Hain program with partial results but no precise modern theorem. Current primary-literature status and redistribution terms remain NEEDS\_REVIEW.
\end{itemize}
\end{document}
